%% Beginning of file 'sample631.tex'
%%
%% Modified 2022 May  
%%
%% This is a sample manuscript marked up using the
%% AASTeX v6.31 LaTeX 2e macros.
%%
%% AASTeX is now based on Alexey Vikhlinin's emulateapj.cls 
%% (Copyright 2000-2015).  See the classfile for details.

%% AASTeX requires revtex4-1.cls and other external packages such as
%% latexsym, graphicx, amssymb, longtable, and epsf.  Note that as of 
%% Oct 2020, APS now uses revtex4.2e for its journals but remember that 
%% AASTeX v6+ still uses v4.1. All of these external packages should 
%% already be present in the modern TeX distributions but not always.
%% For example, revtex4.1 seems to be missing in the linux version of
%% TexLive 2020. One should be able to get all packages from www.ctan.org.
%% In particular, revtex v4.1 can be found at 
%% https://www.ctan.org/pkg/revtex4-1.

%% The first piece of markup in an AASTeX v6.x document is the \documentclass
%% command. LaTeX will ignore any data that comes before this command. The 
%% documentclass can take an optional argument to modify the output style.
%% The command below calls the preprint style which will produce a tightly 
%% typeset, one-column, single-spaced document.  It is the default and thus
%% does not need to be explicitly stated.
%%
%% using aastex version 6.3
\documentclass[linenumbers, twocolumn]{aastex631}

%% The default is a single spaced, 10 point font, single spaced article.
%% There are 5 other style options available via an optional argument. They
%% can be invoked like this:
%%
%% \documentclass[arguments]{aastex631}
%% 
%% where the layout options are:
%%
%%  twocolumn   : two text columns, 10 point font, single spaced article.
%%                This is the most compact and represent the final published
%%                derived PDF copy of the accepted manuscript from the publisher
%%  manuscript  : one text column, 12 point font, double spaced article.
%%  preprint    : one text column, 12 point font, single spaced article.  
%%  preprint2   : two text columns, 12 point font, single spaced article.
%%  modern      : a stylish, single text column, 12 point font, article with
%% 		  wider left and right margins. This uses the Daniel
%% 		  Foreman-Mackey and David Hogg design.
%%  RNAAS       : Supresses an abstract. Originally for RNAAS manuscripts 
%%                but now that abstracts are required this is obsolete for
%%                AAS Journals. Authors might need it for other reasons. DO NOT
%%                use \begin{abstract} and \end{abstract} with this style.
%%
%% Note that you can submit to the AAS Journals in any of these 6 styles.
%%
%% There are other optional arguments one can invoke to allow other stylistic
%% actions. The available options are:
%%
%%   astrosymb    : Loads Astrosymb font and define \astrocommands. 
%%   tighten      : Makes baselineskip slightly smaller, only works with 
%%                  the twocolumn substyle.
%%   times        : uses times font instead of the default
%%   linenumbers  : turn on lineno package.
%%   trackchanges : required to see the revision mark up and print its output
%%   longauthor   : Do not use the more compressed footnote style (default) for 
%%                  the author/collaboration/affiliations. Instead print all
%%                  affiliation information after each name. Creates a much 
%%                  longer author list but may be desirable for short 
%%                  author papers.
%% twocolappendix : make 2 column appendix.
%%   anonymous    : Do not show the authors, affiliations and acknowledgments 
%%                  for dual anonymous review.
%%
%% these can be used in any combination, e.g.
%%
%% \documentclass[twocolumn,linenumbers,trackchanges]{aastex631}
%%
%% AASTeX v6.* now includes \hyperref support. While we have built in specific
%% defaults into the classfile you can manually override them with the
%% \hypersetup command. For example,
%%
%% \hypersetup{linkcolor=red,citecolor=green,filecolor=cyan,urlcolor=magenta}
%%
%% will change the color of the internal links to red, the links to the
%% bibliography to green, the file links to cyan, and the external links to
%% magenta. Additional information on \hyperref options can be found here:
%% https://www.tug.org/applications/hyperref/manual.html#x1-40003
%%
%% Note that in v6.3 "bookmarks" has been changed to "true" in hyperref
%% to improve the accessibility of the compiled pdf file.
%%
%% If you want to create your own macros, you can do so
%% using \newcommand. Your macros should appear before
%% the \begin{document} command.
%%
\newcommand{\vdag}{(v)^\dagger}
\newcommand\aastex{AAS\TeX}
\newcommand\latex{La\TeX}

%% Reintroduced the \received and \accepted commands from AASTeX v5.2
%\received{March 1, 2021}
%\revised{April 1, 2021}
%\accepted{\today}

%% Command to document which AAS Journal the manuscript was submitted to.
%% Adds "Submitted to " the argument.
%\submitjournal{PSJ}

%% For manuscript that include authors in collaborations, AASTeX v6.31
%% builds on the \collaboration command to allow greater freedom to 
%% keep the traditional author+affiliation information but only show
%% subsets. The \collaboration command now must appear AFTER the group
%% of authors in the collaboration and it takes TWO arguments. The last
%% is still the collaboration identifier. The text given in this
%% argument is what will be shown in the manuscript. The first argument
%% is the number of author above the \collaboration command to show with
%% the collaboration text. If there are authors that are not part of any
%% collaboration the \nocollaboration command is used. This command takes
%% one argument which is also the number of authors above to show. A
%% dashed line is shown to indicate no collaboration. This example manuscript
%% shows how these commands work to display specific set of authors 
%% on the front page.
%%
%% For manuscript without any need to use \collaboration the 
%% \AuthorCollaborationLimit command from v6.2 can still be used to 
%% show a subset of authors.
%
%\AuthorCollaborationLimit=2
%
%% will only show Schwarz & Muench on the front page of the manuscript
%% (assuming the \collaboration and \nocollaboration commands are
%% commented out).
%%
%% Note that all of the author will be shown in the published article.
%% This feature is meant to be used prior to acceptance to make the
%% front end of a long author article more manageable. Please do not use
%% this functionality for manuscripts with less than 20 authors. Conversely,
%% please do use this when the number of authors exceeds 40.
%%
%% Use \allauthors at the manuscript end to show the full author list.
%% This command should only be used with \AuthorCollaborationLimit is used.

%% The following command can be used to set the latex table counters.  It
%% is needed in this document because it uses a mix of latex tabular and
%% AASTeX deluxetables.  In general it should not be needed.
%\setcounter{table}{1}

%%%%%%%%%%%%%%%%%%%%%%%%%%%%%%%%%%%%%%%%%%%%%%%%%%%%%%%%%%%%%%%%%%%%%%%%%%%%%%%%
%%
%% The following section outlines numerous optional output that
%% can be displayed in the front matter or as running meta-data.
%%
%% If you wish, you may supply running head information, although
%% this information may be modified by the editorial offices.
%\shorttitle{AASTeX v6.3.1 Sample article}
%\shortauthors{Schwarz et al.}
%%
%% You can add a light gray and diagonal water-mark to the first page 
%% with this command:
%% \watermark{text}
%% where "text", e.g. DRAFT, is the text to appear.  If the text is 
%% long you can control the water-mark size with:
%% \setwatermarkfontsize{dimension}
%% where dimension is any recognized LaTeX dimension, e.g. pt, in, etc.
%%
%%%%%%%%%%%%%%%%%%%%%%%%%%%%%%%%%%%%%%%%%%%%%%%%%%%%%%%%%%%%%%%%%%%%%%%%%%%%%%%%
%\graphicspath{{./}{figures/}}
%% This is the end of the preamble.  Indicate the beginning of the
%% manuscript itself with \begin{document}.

\begin{document}

\title{Color-Magnitude and Hertzsprung-Russell Diagram Morphology of Simulated Star Cluster Formation}

%% LaTeX will automatically break titles if they run longer than
%% one line. However, you may use \\ to force a line break if
%% you desire. In v6.31 you can include a footnote in the title.

%% A significant change from earlier AASTEX versions is in the structure for 
%% calling author and affiliations. The change was necessary to implement 
%% auto-indexing of affiliations which prior was a manual process that could 
%% easily be tedious in large author manuscripts.
%%
%% The \author command is the same as before except it now takes an optional
%% argument which is the 16 digit ORCID. The syntax is:
%% \author[xxxx-xxxx-xxxx-xxxx]{Author Name}
%%
%% This will hyperlink the author name to the author's ORCID page. Note that
%% during compilation, LaTeX will do some limited checking of the format of
%% the ID to make sure it is valid. If the "orcid-ID.png" image file is 
%% present or in the LaTeX pathway, the OrcID icon will appear next to
%% the authors name.
%%
%% Use \affiliation for affiliation information. The old \affil is now aliased
%% to \affiliation. AASTeX v6.31 will automatically index these in the header.
%% When a duplicate is found its index will be the same as its previous entry.
%%
%% Note that \altaffilmark and \altaffiltext have been removed and thus 
%% can not be used to document secondary affiliations. If they are used latex
%% will issue a specific error message and quit. Please use multiple 
%% \affiliation calls for to document more than one affiliation.
%%
%% The new \altaffiliation can be used to indicate some secondary information
%% such as fellowships. This command produces a non-numeric footnote that is
%% set away from the numeric \affiliation footnotes.  NOTE that if an
%% \altaffiliation command is used it must come BEFORE the \affiliation call,
%% right after the \author command, in order to place the footnotes in
%% the proper location.
%%
%% Use \email to set provide email addresses. Each \email will appear on its
%% own line so you can put multiple email address in one \email call. A new
%% \correspondingauthor command is available in V6.31 to identify the
%% corresponding author of the manuscript. It is the author's responsibility
%% to make sure this name is also in the author list.
%%
%% While authors can be grouped inside the same \author and \affiliation
%% commands it is better to have a single author for each. This allows for
%% one to exploit all the new benefits and should make book-keeping easier.
%%
%% If done correctly the peer review system will be able to
%% automatically put the author and affiliation information from the manuscript
%% and save the corresponding author the trouble of entering it by hand.

%\correspondingauthor{August Muench}
%\email{greg.schwarz@aas.org, gus.muench@aas.org}

\author[0009-0004-8200-9910]{Winston Zhang}
\affiliation{University of Virginia, Department of Astronomy \\
Charlottesville, VA 22904, USA}

\author[0000-0002-4725-7589]{Aayush Gautam}
\affiliation{University of Virginia, Department of Astronomy \\
Charlottesville, VA 22904, USA}
\affiliation{Tribhuvan University, Department of Physics \\
Bharatpur 44207, Nepal}

\author[0000-0002-5851-2602]{Juan Farias}
\affiliation{Chalmers University of Technology, Department of Space, Earth \& Environment \\
Gothenburg, Sweden}
\affiliation{University of Texas at Austin, Department of Astronomy \\
Austin, TX 78712, USA}

\author[0000-0003-2874-1196]{Matthew W. Hosek Jr.}
\affiliation{University of California, Los Angeles, Department of Astronomy \\
Los Angeles, CA 90095, USA}

\author[0000-0002-4855-1325]{Joseph J Armstrong}
\affiliation{Chalmers University of Technology, Department of Space, Earth \& Environment \\
Gothenburg, Sweden}

%\author{Aabha Brizawasi}
%\affiliation{Indian Institute of Science Education and Research \\
% Mohali, India}

\author[0000-0002-5306-4089]{Morten Andersen}
\affiliation{European Southern Observatory \\
München, Germany}

\author[0000-0002-3389-9142]{Jonathan C. Tan}
\affiliation{University of Virginia, Department of Astronomy \\
Charlottesville, VA 22904, USA}
\affiliation{Chalmers University of Technology, Department of Space, Earth \& Environment \\
Gothenburg, Sweden}

%% Note that the \and command from previous versions of AASTeX is now
%% depreciated in this version as it is no longer necessary. AASTeX 
%% automatically takes care of all commas and "and"s between authors names.

%% AASTeX 6.31 has the new \collaboration and \nocollaboration commands to
%% provide the collaboration status of a group of authors. These commands 
%% can be used either before or after the list of corresponding authors. The
%% argument for \collaboration is the collaboration identifier. Authors are
%% encouraged to surround collaboration identifiers with ()s. The 
%% \nocollaboration command takes no argument and exists to indicate that
%% the nearby authors are not part of surrounding collaborations.

%% Mark off the abstract in the ``abstract'' environment. 
\begin{abstract}
The morphology of pre--main-sequence stellar populations in color--magnitude diagrams (CMDs) and Hertzsprung--Russell (H--R) diagrams provides a potential diagnostic of the star-formation histories of young stellar clusters. We investigate the relationship between the physical characteristics of a molecular clump and a resultant protocluster's CMD morphology using a forward-modeling framework that combines Nbody6 cluster formation simulations with synthetic photometry generated using SPISEA. We construct synthetic H--R diagrams and JWST/HST CMDs and quantify their vertical spread using a quartile-tail separation statistic.

Analyzing a suite of simulations spanning two cloud surface densities ($\Sigma_{\rm cloud}$) and five star formation efficiencies per free-fall time ($\epsilon_{\rm ff}$), we find that lower $\epsilon_{\rm ff}$ and lower $\Sigma_{\rm cloud}$ produce longer formation timescales, larger intrinsic stellar age spreads, and broader stellar sequences. These trends persist across multiple photometric projections, although the measured spread depends on isochrone geometry. 
% Add the following sentence back if we decide to include M_cl comparison
% A comparison of $300$, $3000$, and $30000\,M_\odot$ protoclusters indicates that increasing cluster mass delays and generally increases the maximum diagram spread. 
We additionally find that unresolved binaries introduce an approximately additive increase in CMD spread without substantially changing its
temporal evolution. At 410 pc, approximately 70--80\% of binary systems remain unresolved under the adopted JWST and HST resolution criteria.

Our results demonstrate that CMD and H--R diagram morphology retains information about cluster formation histories, but diagram width cannot uniquely determine intrinsic age spread without accounting for cloud properties, stellar evolution, and observational resolution.

\end{abstract}

%% Keywords should appear after the \end{abstract} command. 
%% The AAS Journals now uses Unified Astronomy Thesaurus concepts:
%% https://astrothesaurus.org
%% You will be asked to selected these concepts during the submission process
%% but this old "keyword" functionality is maintained in case authors want
%% to include these concepts in their preprints.
\keywords{
Star formation (1569) ---
Star clusters (1567) ---
Pre-main sequence stars (1290) ---
Stellar evolution (1599) ---
Hertzsprung Russell diagram (725) ---
Stellar populations (1622)
}

%% From the front matter, we move on to the body of the paper.
%% Sections are demarcated by \section and \subsection, respectively.
%% Observe the use of the LaTeX \label
%% command after the \subsection to give a symbolic KEY to the
%% subsection for cross-referencing in a \ref command.
%% You can use LaTeX's \ref and \label commands to keep track of
%% cross-references to sections, equations, tables, and figures.
%% That way, if you change the order of any elements, LaTeX will
%% automatically renumber them.
%%
%% We recommend that authors also use the natbib \citep
%% and \citet commands to identify citations.  The citations are
%% tied to the reference list via symbolic KEYs. The KEY corresponds
%% to the KEY in the \bibitem in the reference list below. 


\section{Introduction}
\label{sec:intro}

The formation of stellar clusters is a fundamental component of the star-formation process in galaxies, linking the physical evolution of molecular clouds to the properties of stellar populations. Most stars form in clustered or spatially associated environments, where their evolution is influenced by the dynamical conditions of their parent clouds, as well as interactions with neighboring stars \citep{LadaLada2003,Krumholz2019}. Therefore, understanding how these systems assemble therefore connecting the physical conditions of star-forming molecular clouds to the observable properties of the stellar populations that they produce.

An important question in cluster formation is the timescale over which a stellar population assembles. Star formation may occur relatively rapidly, producing a population whose members have nearly identical ages, or proceed over an extended interval, producing a cluster containing stars with a substantial disparity of formation times. The duration of star formation is expected to depend on the physical properties of the parent molecular cloud, including its density, freefall timescale, and efficiency of converting gas into stars \citep{McKeeOstriker2007,KrumholzTan2007}. Consequently, measurements of stellar age distributions provide a potential means of constraining the physical processes governing cluster formation.

Directly determining the formation history of a young stellar cluster from observations, however, presents several challenges. Individual stellar ages are generally inferred by comparing observed stellar properties with theoretical evolutionary models \citep{Soderblom2010}. This comparison is commonly performed using Hertzsprung--Russell (H--R) diagrams or color--magnitude diagrams (CMDs), in which the observed stellar population is compared with theoretical isochrone curves \citep{Hillenbrand2008}. The location of a star relative to these isochrones provides an estimate of its evolutionary state and, under appropriate assumptions, its age.

The interpretation of these diagrams is complicated by the fact that stars of different ages do not occupy completely distinct regions of parameter space. During pre--main-sequence evolution, changes in stellar luminosity and effective temperature produce age-dependent shifts in the stellar sequence. A population containing stars formed over a range of times may therefore occupy a broader region of an H--R diagram or CMD than a coeval population. In principle, this broadening provides an observational signature of an extended star-formation history \citep{Hillenbrand2008,Reggiani2011}. Nevertheless, the width of a stellar sequence is not determined by age spread alone. The geometry of pre--main-sequence evolutionary tracks changes with age, and the transformation from effective temperature and luminosity into observable photometric magnitudes and colors introduces additional nonlinearities \citep{Bell2012,PecautMamajek2013}. As a result, the same intrinsic age distribution may produce different apparent widths in different CMD projections.

Observational effects introduce further complications. For example, in unresolved binary systems, the combined light of multiple stellar components are attributed to a single photometric source. These systems can appear brighter than individual stars of comparable mass and age, broadening the observed stellar sequence even without an increase in the intrinsic age spread \citep{HurleyTout1998,DaRio2010}. Additionally, effects such as extinction, photometric uncertainty, incompleteness, and background contamination can introduce additional CMD/H--R scatter \citep{DaRio2010,Reggiani2011}. Therefore, distinguishing these contributions from the broadening produced by an extended star-formation history is essential when interpreting observed young clusters.

A useful approach to investigating these effects is forward modeling, in which a physically motivated simulation of cluster formation is transformed into the same observable parameter space used to study real stellar populations \citep{DaRio2010,Masegian2021}. Rather than inferring stellar ages independently from an observed diagram, forward modeling begins with a simulated population whose stellar masses, formation times, and dynamical history are known. Stellar-evolution models can then be used to predict the effective temperatures, luminosities, and photometric magnitudes of the simulated stars at successive stages of cluster evolution.

This approach provides a direct connection between the physical parameters controlling cluster formation and the resulting stellar diagram morphology. It also permits controlled experiments in which individual formation parameters or observational assumptions are varied while other aspects of the stellar population are held fixed. By examining how the synthetic diagrams respond to these changes, one can determine which morphological features are associated with the underlying star-formation history and which arise from the choice of photometric representation or observational treatment.

In this work, we investigate the relationship between protocluster formation histories and synthetic H--R and CMD morphology using numerical simulations of forming stellar clusters. We combine the Nbody6 protocluster simulations of \citet{Farias2019} with stellar-evolution models and synthetic photometry generated using SPISEA \citep{Hosek2020}. Building on the luminosity-spread analysis of \citet{Masegian2021} and the synthetic CMD framework developed by \citet{Zhang2024}, we examine how the vertical extent of the pre--main-sequence stellar population responds to changes in the physical conditions of cluster formation. The principal objective of this study is to establish how the physical parameters and formation histories of simulated protoclusters are reflected in the morphology of their synthetic stellar diagrams.

The remainder of this paper is organized as follows. Section~\ref{sec:related} reviews the numerical and synthetic-photometry studies upon which this work builds. Section~\ref{sec:methods} describes the protocluster simulations, stellar interpolation, synthetic diagram construction, and morphology measurements. Section~\ref{sec:results} presents the evolution of intrinsic stellar age spreads and synthetic diagram widths, their dependence on cluster-formation parameters, and the effects of binary angular resolution. Finally, Section~\ref{sec:conclusion} summarizes the principal findings and discusses their implications for interpreting observed young stellar populations.





\section{Related Work}\label{sec:related}

Previous studies have linked cloud conditions and stellar formation timescales to the age distributions and synthetic diagram morphology of young clusters \citep{Farias2019,Masegian2021,Zhang2024}. The relevant simulation and diagram-analysis methods are discussed in the subsections below.

\section{Methods} \label{sec:methods}

We develop a forward-modeling framework that links protocluster formation simulations to CMDs and H--R diagrams. The methodology combines numerical cluster simulations with stellar-evolution models and synthetic photometry to map simulated stellar masses and ages into theoretical and observable diagram spaces. We then quantify the morphology of the resulting stellar sequences and examine how that morphology depends on the star-formation history of the parent protocluster.

The forward-modeling pipeline combines protocluster simulations from \citet{Farias2019} with synthetic photometry generated using SPISEA \citep{Hosek2020}. Stellar masses and ages extracted from Nbody6 simulations are interpolated onto a grid of theoretical isochrones. The resulting stellar populations are analyzed both in H--R space, using effective temperature and luminosity, and in several photometric CMD projections.

Our primary analysis focuses on a suite of $3000\,M_\odot$ molecular clumps spanning two cloud surface densities and five star formation efficiencies per free-fall time. By analyzing multiple realizations of each parameter combination, we distinguish systematic changes in diagram morphology from stochastic variation between individual simulated clusters. We quantify the vertical spread of the stellar sequence in H--R space and in three CMD projections constructed using JWST and HST photometric filters, and compare these measurements with the intrinsic stellar age distributions of the simulated populations.

We additionally investigate two effects that extend beyond the primary parameter ensemble. First, we compare representative simulations with initial clump masses of $300$, $3000$, and $30000\,M_\odot$ to examine how the characteristic mass scale of a forming cluster influences the development and magnitude of diagram spread. Second, we examine the observational consequences of binary angular resolution by comparing fully resolved, fully unresolved, and telescope-resolution-limited representations of the same simulated population. These experiments allow us to assess both the physical and observational contributions to the measured width of a young stellar sequence.

\subsection{Protocluster Simulations} \label{sec:protocluster-sims}

We utilize the Nbody6 simulations developed by \citet{Farias2019}, which model the gradual formation and early dynamical evolution of embedded star clusters within molecular clouds. These simulations follow the Turbulent Clump Core Accretion (TCCA) framework, based on the Turbulent Core Accretion model of \citet{McKeeTan2003}, in which stars form through core accretion within gravitationally bound, turbulent, and magnetized gas clumps embedded in giant molecular clouds. The framework has subsequently been applied to simulations spanning different cluster masses and environments, including studies of stellar clustering and multiplicity \citep{Gautam2026}.

The parent clumps are modeled as pressure-confined, approximately virialized polytropic spheres, with radial density profiles following $\rho_{\rm cl}(r)\propto r^{-k_\rho}$, where $k_\rho=1.5$. Their characteristic radii, densities, and velocity dispersions depend on the initial clump mass $M_{\rm cl}$ and the surrounding cloud surface density $\Sigma_{\rm cloud}$, the latter of which determines the confining pressure of the clump, such that a higher $\Sigma_{\rm cloud}$ produces a smaller and denser clump at fixed mass. Rather than directly simulating the gas hydrodynamics, the N-body models represent the gravitational influence of the remaining gas through a time-dependent analytical potential, which decreases as star formation progresses.

Stars are gradually introduced according to a prescribed star formation rate, parameterized by the star formation efficiency per free-fall time $\epsilon_{\rm ff}$. The models assume an overall star formation efficiency of $\epsilon=0.5$, with the remaining gas removed as the stellar population forms. The characteristic formation time is therefore

\begin{equation}
    t_\star = \frac{\epsilon}{\epsilon_{\rm ff}}t_{\rm ff,0},
\end{equation}

where $t_{\rm ff,0}$ is the initial free-fall time of the parent clump. Lower $\epsilon_{\rm ff}$ values result in more extended star-formation histories, while higher values produce more rapid formation. Since the initial free-fall time also depends on $M_{\rm cl}$ and $\Sigma_{\rm cloud}$, these parameters jointly determine the duration of cluster formation and the resulting stellar age spread. The connection between cloud properties, formation duration, and intrinsic stellar age spread is demonstrated by \citet{Farias2019}.

The simulations adopt a Kroupa initial mass function and initialize 50\% of stellar systems as primordial binaries. These systems subsequently undergo dynamical evolution, including binary disruption, exchanges, and the formation of new bound pairs through stellar encounters. Thus, the binary population at any given cluster time is not necessarily identical to the primordial population.

Our primary ensemble analysis uses simulations with an initial clump mass of $M_{\rm cl}=3000\,M_\odot$. Additional targeted cluster-mass sensitivity tests use simulations with $M_{\rm cl}=300$, $3000$, and $30000\,M_\odot$, holding $\epsilon_{\rm ff}=0.03$ fixed and considering both $\Sigma_{\rm cloud}=0.1$ and $1.0\,{\rm g\,cm^{-2}}$.

At a specified cluster time $t_{\rm cl}$, an Nbody6 snapshot is converted into a stellar population containing the masses and ages of the stars present at that epoch. These quantities form the input to the stellar-evolution and synthetic-photometry interpolation described below.

\subsection{Simulation Parameter Grid}
\label{sec:simulation-grid}

For the adopted initial clump mass of $M_{\rm cl}=3000\,M_\odot$, the primary simulation suite spans two cloud surface densities,
\begin{equation}
\Sigma_{\rm cloud}
=
0.1,\ 1.0\ {\rm g\,cm^{-2}},
\end{equation}
and five values of the star formation efficiency per free-fall time,
\begin{equation}
\epsilon_{\rm ff}
=
0.01,\ 0.03,\ 0.10,\ 0.30,\ 1.00.
\end{equation}

Each combination of $\Sigma_{\rm cloud}$ and $\epsilon_{\rm ff}$ is represented by 20 independent realizations, identified by seed numbers 00--19. The full analysis therefore contains 200 gradual cluster formation simulations.

For a fixed cloud surface density, simulations with the same seed but different $\epsilon_{\rm ff}$ values contain the same underlying set of stellar masses sampled from the initial mass function. Changing $\epsilon_{\rm ff}$ for a fixed seed primarily changes the times at which these stars are introduced into the cluster. This enables controlled same-seed comparisons in which differences in stellar population morphology can be associated with the star-formation history rather than with a different random realization of the initial mass function.

For ensemble measurements, all 20 realizations of a given $(\Sigma_{\rm cloud},\epsilon_{\rm ff})$ combination are analyzed independently. We report the arithmetic mean across the contributing seeds, while the realization-to-realization scatter is characterized by the sample standard deviation. The resulting standard-deviation bands therefore represent stochastic variation between cluster realizations rather than uncertainty in the ensemble mean.

\begin{deluxetable}{ccc}
\tablecaption{$3000\,M_\odot$ Protocluster Simulation Grid. \textcolor{red}{Refer to Table 1 of same stellar clustering paper and add the relevant columns. Maybe only add $\epsilon_{\rm ff}$=0.03 lines for M300 and M30000 cases since we analyze them.}}
\label{tab:simulation-grid}
\tablehead{
\colhead{$\Sigma_{\rm cloud}$ (${\rm g\,cm^{-2}}$)} &
\colhead{$\epsilon_{\rm ff}$} &
\colhead{Number of realizations}
}
\startdata
0.1 & 0.01 & 20 \\
0.1 & 0.03 & 20 \\
0.1 & 0.10 & 20 \\
0.1 & 0.30 & 20 \\
0.1 & 1.00 & 20 \\
1.0 & 0.01 & 20 \\
1.0 & 0.03 & 20 \\
1.0 & 0.10 & 20 \\
1.0 & 0.30 & 20 \\
1.0 & 1.00 & 20 \\
\enddata
\end{deluxetable}

% The full ensemble analysis presented in this work is based on the $M_{\rm cl}=3000\,M_\odot$ simulation suite. The additional $M_{\rm cl}=300$ and $30000\,M_\odot$ simulations are used only for the controlled mass-dependence experiment described below and are not incorporated into the multi-seed ensemble statistics.

\subsection{Isochrone Grid Construction}

Synthetic stellar populations are generated using SPISEA \citep{Hosek2020}. This follows the synthetic CMD framework of \citet{Zhang2024}, which interpolated simulated stellar masses and ages onto SPISEA isochrones to generate photometric diagrams. The analysis adopts the merged Baraffe--Pisa--Ekstr\"om--PARSEC evolutionary model grid used by the forward-modeling pipeline. Synthetic photometry is calculated with the BT-Settl 2015 stellar-atmosphere models.

The isochrone grid spans ages from 1 to 20 Myr in intervals of 0.5 Myr. We assume solar metallicity, zero extinction ($A_{K_s}=0$), and a distance of 410 pc. The same stellar-evolution and photometric assumptions are held fixed throughout the simulation ensemble so that changes in diagram morphology arise from changes in the underlying simulated stellar population rather than from changes in the synthetic-observation model.

Photometry is generated in the JWST F070W, F182M, and F200W filters and the HST F555W and F814W filters. 

\subsection{Photometric Interpolation}

For each simulated star with mass $m$ and age $a$, stellar parameters and photometric magnitudes are obtained by bilinear interpolation within the isochrone grid. For a 1--20 Myr, 0.01--120 $M_\odot$ age--mass pair:

\begin{enumerate}
    \item identify the two grid isochrones whose ages bracket $a$;
    \item interpolate in stellar mass along each bracketing isochrone;
    \item interpolate in age between the two mass-interpolated points.
\end{enumerate}

The resulting stellar properties include effective temperature, luminosity, surface gravity, and magnitudes in all selected filters. Stars younger than the minimum age represented by the stellar-evolution grid are evaluated at the youngest available grid age, while stars outside the supported upper age range are excluded.

The interpolation procedure is illustrated schematically in Figure~\ref{fig:iso-interp-example}.

\begin{figure}[ht]
    \centering
    \includegraphics[width=1\linewidth]{science-images/fig_interpolation_schematic_f162m_f182m.png}
    \caption{Example bilinear interpolation using neighboring stellar isochrones to obtain the synthetic properties of a star whose mass and age do not coincide exactly with the discrete isochrone grid.}
    \label{fig:iso-interp-example}
\end{figure}





\subsection{Treatment of Binary Systems}
\label{sec:binaries}
The Nbody6 simulations include binary systems whose observational appearance depends on whether the individual stellar components can be spatially resolved. We therefore consider both resolved and unresolved representations of binary systems when constructing synthetic CMDs.

\paragraph{Resolved binaries:}
For a resolved binary, the primary and companion are treated as independent stellar sources. Each component is interpolated separately according to its stellar mass and age and contributes an independent point to the CMD.

\paragraph{Unresolved binaries:}
For an unresolved binary, the primary and companion are represented as a single photometric source. Each component is interpolated separately, then the components have their magnitude brightnesses independently converted to fluxes, which are then summed in each photometric band before converting back to the magnitude scale. For a given filter,

\begin{equation}
F_{\rm binary}=F_1+F_2,
\end{equation}

or equivalently,

\begin{equation}
m_{\rm binary}
=
-2.5\log_{10}
\left(
10^{-0.4m_1}+10^{-0.4m_2}
\right),
\end{equation}

apart from the common photometric zero point.

The resulting combined magnitudes are then used to construct the CMD color. Because the combined source is brighter than either component individually, unresolved multiplicity can broaden the stellar sequence and alter the measured quartile-tail spread, as explained in subsequent subsections.

Figure~\ref{fig:cmd-binaries} illustrates the limiting cases in which all binary systems are treated as resolved or all are treated as unresolved.

\begin{figure*}
    \centering
    \includegraphics[width=1\linewidth]{science-images/fig_cmd_binaries_f162m_f182m.png}
    \caption{Example synthetic CMDs demonstrating the difference between treating binary systems as resolved individual stellar components and unresolved combined photometric sources. In the unresolved case, the primary and companion fluxes are summed before their combined magnitude and color are calculated. The presence of unresolved binary systems, due to their nature as being a combined light source, results in the apparent effect of a wider spread along the vertical axis.}
    \label{fig:cmd-binaries}
\end{figure*}


\subsubsection{Angular-Resolution Treatment}
\label{sec:binary-resolution-method}
The all-resolved and all-unresolved cases provide useful limiting examples, but a real observation generally lies between these extremes. We therefore perform an additional angular-resolution experiment to estimate which binary systems would be resolved by JWST or HST for a nearby young cluster.

\textbf{
This experiment uses a representative simulation with}

\begin{equation}
M_{\rm cl}=3000\,M_\odot,
\qquad
\Sigma_{\rm cloud}=0.1\,{\rm g\,cm^{-2}},
\qquad
\epsilon_{\rm ff}=0.03,
\end{equation}

\textbf{using seed 00 \textcolor{red}{(this sentence will probably change after redoing the $M_{cl}$ experiment)}.} The experiment is intended as a test of the observational effect of binary resolution rather than as a second full ensemble analysis.

The binary population is determined independently at each simulation snapshot. Thus, the calculation follows the current binary systems rather than assuming that all primordial pairs remain unchanged throughout the simulation. Binary disruption, dynamical formation, exchanges, and orbital evolution can therefore alter both the binary population and the instantaneous component separations.

For each current binary, the projected separation is calculated in three orthogonal viewing directions. We refer to these as the XY, XZ, and YZ sky-plane projections. A binary is classified as unresolved when its projected separation is smaller than the angular resolution defined by the Rayleigh criterion,

 
\begin{equation}
\theta_{\rm R} = 1.22\frac{\lambda}{D},
\end{equation}

where $\lambda$ is the adopted observing wavelength and $D$ is the telescope primary-mirror diameter.

The cluster is placed at a distance of 410 pc, approximately corresponding to the distance adopted for the Orion Nebula Cluster comparison. The angular Rayleigh limit is converted into a physical projected-separation threshold at this distance. 

For each CMD, the longest-wavelength filter entering the diagram is used to define the resolution limit. The JWST
$(F070W-F200W)$ versus F200W and
$(F182M-F200W)$ versus F200W diagrams therefore use the F200W wavelength and a telescope diameter of 6.5 m. The HST $(F555W-F814W)$ versus F814W diagram uses F814W and a telescope diameter of 2.4 m.

Five binary-resolution treatments are compared at each cluster time: the all-resolved limit, the all-unresolved limit, and Rayleigh-limited observations in the XY, XZ, and YZ projections. For a pair classified as unresolved, the two stellar components are combined in flux independently in every photometric band before the CMD color is calculated. Resolved systems contribute their two stellar components independently.

The same fixed CMD measurement bins and quartile-tail spread statistic described in Section~\ref{sec:spread-metric} are then applied to all five resolution treatments. The experiment is restricted to the three CMDs because spatial resolving power is an observational imaging effect rather than an intrinsic property of H--R space.

Our usage of the Rayleigh criterion to only consider bound binaries has the caveat that we neglect to consider gravitationally unbound, apparent unresolved binaries. Doing so would involve devising a searching algorithm that would draw annuli around each star in the sky plane projection according to the Rayleigh threshold. This approach would account for all orders of unresolved multiplicity (i.e., binaries, triples, and above); however, doing so would require convolving the stellar population with the point-spread function of the longest filter considered in the photometric analysis. Distant background stars would also need to be considered alongside the PSF convolution. Since the focus of this paper is not on synthetic image generation, and the simulations used only include binaries, these edge-case considerations are largely beyond the scope of this work.





\subsection{Synthetic Diagram Construction}

For every interpolated stellar population, we construct both theoretical H--R diagrams and synthetic CMDs. The H--R representation uses $T_{\rm eff}$ versus $\log_{10}(L/L_\odot)$,
while photometric colors are formed from differences between synthetic magnitudes.

The complete ensemble analysis considers the following three CMD spaces: 

\begin{enumerate}
    \item (F070W-F200W) versus F200W, 
    \item (F182M-F200W) versus F200W,
    \item (F555W-F814W) versus F814W.
\end{enumerate}

Together with the H--R diagram, these projections allow us to determine whether trends present in physical stellar-parameter space remain visible after transformation into observable photometric space.

\subsection{Definition of the Measurement Region}
\label{sec:measurement-region}

The morphology statistic is measured over a restricted portion of the low-mass pre--main-sequence population. A star is eligible for the diagram-spread analysis only when 3000 K $\leq T_{ eff} \leq$ 3600 K. The measurement geometry is fixed for all cluster ages and all simulations. This is motivated by \citet{Masegian2021}, which adopts the same measurement region and morphology statistic as a method of comparison to the Orion Nebula cluster, whose initial mass function peaks in the mass range corresponding to this range of effective temperature \citep{DaRio2016}. 

In H--R space, the interval from 3000 to 3600 K is divided directly into six 100 K bins,
\begin{equation}
3000,\ 3100,\ \ldots,\ 3600\ {\rm K}.
\end{equation}

The corresponding CMD binning is defined differently because color is a nonlinear function of effective temperature. For each CMD projection, the 1 Myr reference isochrone is used to determine the color corresponding to $T_{\rm eff}=3000$ K and the color corresponding to $T_{\rm eff}=3600$ K. The fixed color interval between these two values is then divided into six equal-width bins. These color boundaries are held fixed for every subsequent simulation snapshot and cluster age. Individual stars are assigned to the bins according to their actual interpolated color.

This procedure provides a common low-mass stellar interval across the different diagram spaces while avoiding age-dependent movement of the bin boundaries. Figure~\ref{fig:analysis-geometry} illustrates the isochrone geometry, representative stellar population, and fixed measurement boundaries used in each diagram.

\begin{figure*}
    \centering
    \includegraphics[width=1\linewidth]{science-images/fig_analysis_space_isochrone_geometry.png}
    \caption{Isochrone geometry, fixed measurement bins, and a representative simulated stellar population in the four diagram spaces used for the ensemble analysis. Isochrones are colored by age. In H--R space the boundaries correspond to six 100 K intervals between 3000 and 3600 K. In each CMD, the boundaries are six equal-width color bins defined between the colors corresponding to 3000 and 3600 K on the 1 Myr reference isochrone.}
    \label{fig:analysis-geometry}
\end{figure*}

\subsection{Dependence on Diagram Space \textbf{\textcolor{blue}{(possibly merge this with previous subsection)}}}

The same stellar population can have substantially different apparent morphology depending on the set of filters in which the color-magnitude diagrams are made. Figure~\ref{fig:analysis-geometry} compares the H--R diagram to the three CMD projections included in the ensemble analysis.

In H--R space, the low-mass pre--main-sequence tracks remain relatively well ordered in effective temperature and luminosity. The transformation into photometric space alters both the slope and curvature of these tracks. In particular, the same $3000$--$3600$ K stellar interval can span very different ranges in color depending on the adopted filter pair.

This motivates the use of a separately defined fixed color interval for each CMD rather than treating equal color intervals in different filters as physically equivalent. The $(F182M-F200W)$ versus F200W projection provides a comparatively well-behaved mapping of the selected temperature range and remains a useful principal photometric representation, while the additional JWST and HST projections provide tests of how robust the inferred trends are to the choice of observable space.

\subsection{Diagram Morphology Quantification}
\label{sec:spread-metric}

We quantify the vertical extent of the stellar sequence using the separation between the lower and upper quartiles of the distribution in each horizontal bin. Motivated by \citet{Masegian2021}, who linked luminosity spread in synthetic H--R diagrams to intrinsic stellar age spread, we adapt this quartile-tail approach for both H--R and CMD space.

For a valid horizontal bin $j$, the 25th and 75th percentiles of the vertical coordinate are first calculated. We then compute the mean vertical coordinate of the stars in the lower quartile and the mean vertical coordinate of the stars in the upper quartile.

For a CMD, where the vertical coordinate is magnitude, the bin spread is
\begin{equation}
S_{m,j}
=
\left\langle m\right\rangle_{{\rm faint},j}
-
\left\langle m\right\rangle_{{\rm bright},j},
\end{equation}
such that a broader stellar sequence gives a larger positive value.

For an H--R diagram, the analogous luminosity spread is
\begin{equation}
S_{L,j}
=
\left\langle \log(L/L_\odot)\right\rangle_{{\rm upper},j}
-
\left\langle \log(L/L_\odot)\right\rangle_{{\rm lower},j}.
\end{equation}

A bin must contain at least eight eligible stars to contribute to the measurement. The overall spread is the arithmetic mean over all valid bins,
\begin{equation}
S
=
\frac{1}{N_{\rm valid}}
\sum_{j=1}^{N_{\rm valid}}S_j.
\end{equation}

Using the means of the outer quartiles suppresses sensitivity to individual outliers while retaining sensitivity to broadening of the stellar sequence. The fixed six-bin construction also ensures that the same portion of diagram space is sampled consistently across simulations and cluster ages.

\subsection{Intrinsic Stellar Age Spread}

The intrinsic stellar age spread of each finite-formation simulation snapshot is quantified independently of the diagram morphology as
\begin{equation}
\Delta t_{90}
=
t_{95}-t_5,
\end{equation}
where $t_{95}$ and $t_5$ are the 95th and 5th percentiles of the stellar-age distribution.

The age percentiles are calculated from the full successfully interpolated stellar population in the snapshot, rather than only from the 3000--3600 K subset used for the diagram-spread measurement. This definition suppresses sensitivity to a small number of extreme-age stars while characterizing the duration represented by the bulk of the stellar population.

\subsection{Ensemble Analysis}

The spread measurement is evaluated at cluster times from 1 to 20 Myr in 0.5 Myr intervals for every available realization, yielding a time-dependent measurement for each $(\Sigma_{\rm cloud},\epsilon_{\rm ff},{\rm seed})$ combination and diagram space.

For each $(\Sigma_{\rm cloud},\epsilon_{\rm ff})$ family, the measurements from the 20 seeds are combined to determine the ensemble mean and sample standard deviation at each cluster time. We compare these ensemble measurements in three complementary ways:

\begin{enumerate}
    \item the evolution of the intrinsic stellar age spread $\Delta t_{90}$ with cluster time;
    \item the evolution of diagram spread with cluster time;
    \item the direct relation between diagram spread and intrinsic stellar age spread.
\end{enumerate}

This ensemble treatment separates systematic trends with the cloud parameters from stochastic differences between individual realizations and allows the dependence of diagram morphology on $\epsilon_{\rm ff}$ and $\Sigma_{\rm cloud}$ to be evaluated across the complete simulation grid.

\subsection{Cluster-Mass Sensitivity Test}
\label{sec:mass-test}

In addition to the primary $M_{\rm cl}=3000\,M_\odot$ ensemble analysis, we perform a targeted experiment to isolate the effect of initial cluster mass on the evolution of diagram spread. We compare simulations with

\begin{equation}
    M_{\rm cl} = 300,\ 3000,\ 30000\,M_\odot
\end{equation}

while fixing the star formation efficiency per free-fall time to $\epsilon_{\rm ff}=0.03$.

The comparison is performed independently for $\Sigma_{\rm cloud}=0.1$ and $1.0\,{\rm g\,cm^{-2}}$.

For each mass and surface density, one representative simulation realization is analyzed using the same stellar interpolation, diagram definitions, measurement region, and quartile-tail spread statistic used in the primary ensemble analysis. The purpose of this experiment is therefore not to measure the ensemble distribution at each cluster mass, but to determine whether a first-order systematic change in diagram morphology occurs when the mass scale of the forming cluster is varied while the other principal formation parameters are held fixed.

A single realization is used at each mass because the number of available simulations differs substantially between the three mass suites. For each $(\Sigma_{\rm cloud},\epsilon_{\rm ff})$ combination, the $300\,M_\odot$ suite contains 200 realizations, the $3000\,M_\odot$ suite contains 20, and the $30000\,M_\odot$ suite contains only two. An all-realization comparison would therefore have strongly unequal sampling precision between the three mass scales. Conversely, restricting every mass to only two realizations would discard most of the available sampling at lower masses while still providing only a weak estimate of the realization variance at $30000\,M_\odot$.

We consequently interpret the cluster-mass comparison as a qualitative sensitivity test rather than an ensemble measurement and do not attach realization-to-realization uncertainty bands to these curves. This treatment is also motivated by the relatively modest seed-to-seed scatter measured in the $3000\,M_\odot$ ensemble compared with the systematic differences observed between the three cluster masses.





\section{Results \& Analysis} \label{sec:results}

The forward-modeling framework successfully maps the simulated stellar populations into both H--R and color--magnitude space. The resulting diagrams evolve systematically with the stellar age distribution and retain measurable signatures of the cloud parameters that regulate cluster formation.

\subsection{Synthetic Cluster Morphology}

Figure~\ref{fig:representative-cluster} shows a representative simulated cluster in the $(F182M-F200W)$ versus F200W CMD together with the synthetic isochrone grid. Individual simulated stars populate a broad pre--main-sequence locus bounded by isochrones of different ages. The population does not lie on a single-age sequence because star formation proceeds over a finite interval in the simulations.

\begin{figure}
    \centering
    \includegraphics[width=1.0\linewidth]{science-images/fig_representative_cluster_f182m_f200w.png}
    \caption{Representative synthetic cluster in $(F182M-F200W)$ versus F200W space, overlaid on the 1--20 Myr isochrone grid. Isochrones are colored by age using the same color scale adopted throughout the diagram-geometry figures.}
    \label{fig:representative-cluster}
\end{figure}

The morphology of an individual cluster also changes with cluster time. Figure~\ref{fig:cmd-age-evolution} follows one realization at six epochs in the $(F182M-F200W)$ versus F200W projection. The same isochrone grid is shown in each panel, allowing the evolution of the population to be separated visually from changes in the underlying reference tracks.

At early times, the population contains a substantial fraction of very young, luminous pre--main-sequence stars. As the cluster ages, the entire sequence evolves toward fainter magnitudes while the stellar age distribution also changes as star formation proceeds and ultimately subsides. The resulting CMD width therefore reflects both stellar evolution and the duration of star formation.

\begin{figure*}
    \centering
    \includegraphics[width=1\linewidth]{science-images/fig_cluster_time_evolution_f182m_f200w.png}
    \caption{Evolution of one simulated cluster realization in $(F182M-F200W)$ versus F200W space. Each panel shows the same 1--20 Myr isochrone grid and the interpolated stellar population at a different cluster time.}
    \label{fig:cmd-age-evolution}
\end{figure*}

\subsection{Evolution of the Intrinsic Stellar Age Spread}

Figure~\ref{fig:age-spread-vs-time} shows the ensemble-averaged intrinsic age spread as a function of cluster time for both cloud surface densities and all five values of $\epsilon_{\rm ff}$. Each curve represents the mean across the available seeds, and the shaded region represents the sample standard deviation among realizations.

A clear ordering with $\epsilon_{\rm ff}$ emerges. Simulations with lower $\epsilon_{\rm ff}$ develop larger stellar age spreads because star formation proceeds over a longer interval. Conversely, the highest-efficiency simulations form their stellar populations more rapidly and consequently maintain substantially narrower age distributions.

The cloud surface density introduces a second systematic dependence. At fixed $\epsilon_{\rm ff}$, the $\Sigma_{\rm cloud}=0.1\,{\rm g\,cm^{-2}}$ simulations generally develop larger age spreads than the $\Sigma_{\rm cloud}=1.0\,{\rm g\,cm^{-2}}$ simulations. The lower-density clouds therefore produce more temporally extended stellar populations in the simulation suite considered here. 
\textcolor{red}{Once we present the details of our Turbulent Clump Core Accretion models in Methods, it will be clearer to the reader that for the same $M_{\rm cl}$ and $\epsilon_{\rm ff}$, $\Sigma_{\rm cloud}$ has shorter free-fall timescales ($t_{\rm ff}$) and thus shorter star formation time ($t_\star = (\epsilon/\epsilon_{\rm ff})t_{\rm ff}$). This is the reason why age spreads are larger for smaller $\Sigma_{\rm cloud}$.} 
\textbf{\textcolor{blue}{This effect is demonstrated in \citet{Farias2019}, and transitively has further effects on the CMD.}}
 
At early cluster times, $\Delta t_{90}$ grows as newly formed stars continue to extend the age distribution. Once star formation has largely ceased, the difference between the oldest and youngest major components of the population approaches a characteristic value and the age-spread curves flatten.

\begin{figure*}
    \centering
    \includegraphics[width=1\linewidth]{science-images/fig_ensemble_age_spread_vs_time.png}
    \caption{Ensemble intrinsic stellar age spread $\Delta t_{90}$ as a function of cluster time. Separate panels show $\Sigma_{\rm cloud}=0.1$ and $1.0\,{\rm g\,cm^{-2}}$. Curves show the mean across the 20 simulation realizations and shaded regions show the sample standard deviation among seeds. Lower $\epsilon_{\rm ff}$ produces systematically larger age spreads, while the lower cloud-surface-density simulations also exhibit larger age spreads at fixed $\epsilon_{\rm ff}$.}
    \label{fig:age-spread-vs-time}
\end{figure*}

\subsection{Diagram Spread as a Function of Cluster Time}

The same hierarchy appears in the morphology of the synthetic stellar diagrams. Figure~\ref{fig:spread-vs-time} shows the ensemble quartile-tail spread as a function of cluster time in H--R space and in all three CMD projections.

Across the diagram spaces, lower values of $\epsilon_{\rm ff}$ generally produce larger measured spreads. This mirrors the intrinsic age-spread hierarchy and demonstrates that the extended star-formation histories of the low-efficiency models remain encoded in the morphology of the stellar sequence after transformation into photometric space.

Cloud surface density produces a corresponding systematic effect. For a fixed $\epsilon_{\rm ff}$, clusters formed in the lower-surface-density $\Sigma_{\rm cloud}=0.1\,{\rm g\,cm^{-2}}$ models generally exhibit larger diagram spreads than their $\Sigma_{\rm cloud}=1.0\,{\rm g\,cm^{-2}}$ counterparts. Thus the ordering identified directly in the stellar-age distributions is also recovered from the H--R and CMD morphology.

The absolute value and detailed temporal behavior of the spread depend on diagram space. This is expected because the stellar-evolution tracks are transformed nonlinearly into different photometric coordinates. The agreement in the relative ordering of the simulation families is therefore more physically informative than equality of the numerical spread measured in different diagrams.

Some temporal structure also remains after averaging over 20 realizations. Diagram spread is not simply a monotonic function of cluster time because the morphology is affected simultaneously by the development of the age distribution and by the evolution of the pre--main-sequence isochrones themselves. Nevertheless, the ensemble analysis reveals a persistent separation between simulation families with different formation histories.

\begin{figure*}
    \centering
    \includegraphics[width=1\linewidth]{science-images/fig_ensemble_spread_vs_time_all_spaces.png}
    \caption{Ensemble quartile-tail diagram spread as a function of cluster time in H--R space and the three analyzed CMD projections. Columns correspond to the two cloud surface densities. Gradual formation models show the mean across simulation seeds, with shaded regions indicating the sample standard deviation. Lower $\epsilon_{\rm ff}$ and lower $\Sigma_{\rm cloud}$ produce systematically broader stellar sequences.}
    \label{fig:spread-vs-time}
\end{figure*}

\subsection{Diagram Spread as a Tracer of Intrinsic Age Spread}

The connection between the observable morphology and the underlying stellar-age distribution is shown more directly in Figure~\ref{fig:spread-vs-age-spread}. Rather than plotting both quantities against cluster time, the ensemble mean diagram spread is plotted directly against the corresponding ensemble mean $\Delta t_{90}$.

In H--R space, increasing intrinsic age spread is accompanied by increasing luminosity spread over much of the simulation sequence. This reproduces the qualitative behavior identified by \citet{Masegian2021}: a population containing stars formed over a longer interval occupies a broader vertical region of the pre--main-sequence H--R diagram.

The same relationship remains visible after transformation into CMD space. In particular, the $(F182M-F200W)$ versus F200W projection retains a clear connection between stellar age spread and photometric sequence width. The F070W--F200W and HST F555W--F814W projections show that the detailed shape and normalization of the relation depend on filter choice, but the underlying tendency for larger intrinsic age spreads to produce broader diagrams remains present.

The simulation families occupy different regions of this relation because both $\epsilon_{\rm ff}$ and $\Sigma_{\rm cloud}$ regulate the duration of star formation. Low-$\epsilon_{\rm ff}$ models preferentially occupy the high-$\Delta t_{90}$, high-spread portion of the diagrams, whereas high-$\epsilon_{\rm ff}$ models are concentrated toward smaller age spreads and narrower sequences. Likewise, the lower-density cloud models are displaced toward larger age and diagram spreads relative to the higher-density models.

\begin{figure*}
    \centering
    \includegraphics[width=1\linewidth]{science-images/fig_ensemble_spread_vs_age_spread_all_spaces.png}
    \caption{Ensemble diagram spread as a function of intrinsic stellar age spread $\Delta t_{90}$ for the H--R diagram and three CMD projections. Each sequence traces the evolution of a simulation family through cluster time. Larger intrinsic age spreads generally correspond to broader stellar sequences, with low-$\epsilon_{\rm ff}$ and low-$\Sigma_{\rm cloud}$ models preferentially occupying the high-spread region. Error bars indicate the realization-to-realization sample standard deviation. \textcolor{red}{AG: For this figure, you could try log-scale on the x-axis to check if that shows the high $\epsilon_{\rm ff}$ models better. They will stay as vertical lines, however the value of age-spread could become clearer to see.}}
    \label{fig:spread-vs-age-spread}
\end{figure*}

\subsection{Effect of Binary Angular Resolution}
\label{sec:binary-resolution-results}
The limiting binary treatments shown in Figure~\ref{fig:cmd-binaries} demonstrate qualitatively that unresolved systems broaden a synthetic CMD. The angular-resolution experiment allows this effect to be quantified using a telescope-dependent observational criterion.

Figure~\ref{fig:binary-resolution-summary} compares the evolution of the quartile-tail spread for the all-resolved and all-unresolved limits with the three Rayleigh-limited viewing projections. The same comparison is shown for all three CMD spaces considered in the experiment.

Treating all binaries as unresolved systematically increases the measured CMD spread relative to the all-resolved case. However, the overall time dependence of the spread remains very similar between the two limiting treatments. Unresolved multiplicity therefore behaves primarily as an approximately additive offset in the spread metric rather than producing a qualitatively different temporal evolution. In other words, the underlying evolution of the stellar age distribution remains visible, but the absolute normalization of the measured CMD spread is shifted upward when binary components are combined.

The realistic Rayleigh-limited cases lie substantially closer to the all-unresolved limit than to the all-resolved limit. For the representative simulation at 410 pc, approximately 70--80\% of the current binary population is unresolved over much of the analyzed interval. Because most binaries therefore contribute as combined photometric sources, the resulting CMD spread closely approaches the all-unresolved extreme.

The binary-resolution result is also largely insensitive to the viewing orientation of the cluster. The XY, XZ, and YZ projections produce similar unresolved fractions and nearly overlapping spread curves. Individual binary systems can change between resolved and unresolved classifications as their projected separation changes with viewing direction and orbital evolution, but these differences have little effect on the aggregate CMD-spread measurement.

This behavior is present in all three tested photometric projections:
$(F070W-F200W)$ versus F200W,
$(F182M-F200W)$ versus F200W, and
$(F555W-F814W)$ versus F814W.
Although the numerical value of the spread and the instrumental resolution threshold differ between filter combinations, the same qualitative result is recovered in each case: unresolved binaries increase the spread, realistic angular resolution places the observation closer to the all-unresolved limit, and changing the viewing axis produces only a small effect.

\begin{figure*}
    \centering
    \includegraphics[width=1\linewidth]
    {science-images/fig_binary_resolution_summary.png}
    \caption{Effect of binary angular resolution on the measured CMD spread for the representative
    $M_{\rm cl}=3000\,M_\odot$,
    $\Sigma_{\rm cloud}=0.1\,{\rm g\,cm^{-2}}$,
    $\epsilon_{\rm ff}=0.03$, seed 00 simulation.
    The top row shows the quartile-tail spread as a function of cluster time for the all-resolved and all-unresolved limiting treatments and for Rayleigh-limited observations in the XY, XZ, and YZ sky-plane projections. The bottom row shows the corresponding fraction of current binary systems classified as unresolved. Columns correspond to
    $(F070W-F200W)$ versus F200W,
    $(F182M-F200W)$ versus F200W, and
    $(F555W-F814W)$ versus F814W.
    Approximately 70--80\% of binary systems are unresolved for the representative cluster at 410 pc, causing the Rayleigh-limited spread to lie substantially closer to the all-unresolved limit than to the all-resolved limit. The three viewing projections produce similar results, indicating little dependence on viewing orientation. \textbf{\textcolor{red}{INCREASE FONT SIZE ON THIS FIGURE}}}
    \label{fig:binary-resolution-summary}
\end{figure*}

The effect of these different binary treatments on the appearance of the stellar population is illustrated directly in Figure~\ref{fig:binary-resolution-time-evolution}. The figure follows the same representative simulation through six cluster times in $(F182M-F200W)$ versus F200W space. Each panel contains the same 1--20 Myr isochrone grid, allowing differences between the three binary-resolution treatments to be compared against the underlying stellar-evolution sequence.

In the all-resolved case, every stellar component appears as an independent source. In the all-unresolved case, the two components of each binary system are replaced by a single source whose flux is the sum of the component fluxes. The resulting systems are displaced toward brighter F200W magnitudes, visually producing the broader vertical distribution reflected in the spread metric.

The Rayleigh-limited XY projection produces an intermediate population, but its morphology more closely resembles the all-unresolved case than the all-resolved case. Most primordial binary systems are represented by the dark-blue unresolved population, while the smaller subset resolved by the JWST Rayleigh criterion appears as light-blue individual components. Dynamically formed binaries are shown separately in dark red regardless of whether they are resolved or unresolved. The relative scarcity of light-blue sources compared with dark-blue sources provides a direct visual counterpart to the high unresolved fractions measured in Figure~\ref{fig:binary-resolution-summary}.

The time sequence also illustrates that binary resolution changes primarily the width and normalization of the stellar sequence rather than its overall evolution. All three treatments follow the same progression through the isochrone grid as the cluster ages, while unresolved flux combination adds a persistent population displaced toward brighter magnitudes. This visual behavior is consistent with the approximately additive shift between the resolved and unresolved spread curves.

\begin{figure*}
    \centering
    \includegraphics[width=1.0\linewidth]
    {science-images/fig_binary_resolution_cmd_time_evolution_f182m_f200w.png}
    \caption{Abridged time evolution of the representative cluster in
    $(F182M-F200W)$ versus F200W under three treatments of binary angular
    resolution. Rows correspond to cluster times of 1, 5, and
    10 Myr, and the same 1--20 Myr isochrone grid is shown in each panel.
    The left column treats all binary components as resolved and plots all
    stellar sources in black. The middle column treats all binary systems as
    unresolved; single stars are shown in black, unresolved primordial
    binaries in dark blue, and dynamically formed binaries in dark red.
    The right column applies the JWST Rayleigh criterion using the XY
    sky-plane projection. In this case, single stars are black, resolved
    primordial binary components are light blue, unresolved primordial
    binaries are dark blue, and dynamically formed binaries are dark red
    regardless of their resolution state. The Rayleigh-limited population
    visually resembles the all-unresolved limit because most binary systems
    remain unresolved at the adopted distance of 410 pc. Also note the seldom appearances of dynamically formed binaries from snapshot to snapshot, which shows how rare and short-lived such systems are in these cluster simulations. This is consistent with \citet{Farias2023}, which finds that less than 1\% of binaries are dynamically formed in $3000 M_\odot$ cluster simulations. \textbf{\textcolor{red}{INCREASE LEGEND FONT SIZE ON THIS FIGURE}}}
    \label{fig:binary-resolution-time-evolution}
\end{figure*}

Together, Figures~\ref{fig:binary-resolution-summary} and \ref{fig:binary-resolution-time-evolution} show that unresolved multiplicity is an important systematic when interpreting the absolute value of a CMD-spread measurement. The quantitative spread curves demonstrate that binary blending produces an approximately additive offset, while the time-evolution CMDs show the stellar populations responsible for that offset. The physical trend with cluster evolution remains visible, but the observed sequence is broadened as binary components are combined into unresolved photometric sources. Consequently, synthetic and observed clusters should be compared using consistent assumptions about instrumental angular resolution and binary blending.

\subsection{Dependence on Star Formation Efficiency and Cloud Surface Density}
The ensemble results demonstrate that the diagram spread is sensitive to both principal cloud parameters explored in the simulation suite.

The strongest systematic ordering is associated with $\epsilon_{\rm ff}$. Lower star formation efficiencies per free-fall time correspond to more prolonged formation histories. These simulations therefore contain stars spanning a larger range of formation times and develop larger values of $\Delta t_{90}$. The increased age diversity broadens the pre--main-sequence population vertically in H--R space and produces a corresponding increase in magnitude spread in the synthetic CMDs.

Cloud surface density modifies the same relationship. At otherwise fixed $\epsilon_{\rm ff}$, the lower-density $\Sigma_{\rm cloud}=0.1\,{\rm g\,cm^{-2}}$ simulations exhibit larger intrinsic age spreads and larger diagram spreads than the $\Sigma_{\rm cloud}=1.0\,{\rm g\,cm^{-2}}$ simulations. The synthetic diagrams therefore retain information not only about the duration of star formation as controlled by $\epsilon_{\rm ff}$, but also about the physical conditions of the parent cloud.

These trends are recovered after averaging over the full 20-seed ensembles and are therefore not attributable to a particular stochastic sampling of the stellar initial mass function. Individual realizations exhibit significant scatter, particularly at some cluster ages and in some photometric projections, but the ordering of the ensemble means remains systematic.

\subsection{Dependence on Cluster Mass}
\label{sec:mass-results}

The primary ensemble analysis holds the initial cluster mass fixed at $M_{\rm cl}=3000\,M_\odot$. To determine whether the characteristic spread also depends on the mass scale of the forming system, we repeat the time-dependent spread analysis for representative $M_{\rm cl}=300$, $3000$, and $30000\,M_\odot$ simulations at fixed $\epsilon_{\rm ff}=0.03$. Figure~\ref{fig:mass-spread-evolution} shows the result for both cloud surface densities in H--R space and in the three CMD projections.

A systematic dependence on cluster mass is visible in both surface-density environments. Increasing $M_{\rm cl}$ shifts the development of the diagram spread toward later cluster times and generally allows the population to reach a larger maximum spread. The $300\,M_\odot$ simulations develop and reach their characteristic maximum spread earliest, the $3000\,M_\odot$ models evolve on an intermediate timescale, and the $30000\,M_\odot$ simulations continue building a broad stellar population for substantially longer. \textbf{\textcolor{blue}{This is due to the star cluster formation timescale ($t_\star$) scaling with $M_{\rm cl}$ in the Turbulent Clump Core Accretion models \citep{Farias2019}.}}

This behavior follows naturally from the formation timescales of the underlying simulation models. At fixed $\Sigma_{\rm cloud}$ and $\epsilon_{\rm ff}$, increasing the total clump mass increases the characteristic time required to assemble the stellar population. Star formation therefore remains active for a longer absolute time in the more massive systems. Because stars continue to be introduced over this extended interval, the stellar population develops a larger absolute age spread before formation terminates, which in turn produces a broader vertical distribution in both H--R and CMD space.

The mass dependence is therefore conceptually distinct from the
$\epsilon_{\rm ff}$ dependence even though both act through the duration of star formation. Lowering $\epsilon_{\rm ff}$ prolongs star formation by reducing the rate at which gas is converted into stars per free-fall time, whereas increasing $M_{\rm cl}$ lengthens the characteristic physical formation timescale of the system at otherwise fixed parameters. Both effects can consequently increase the absolute stellar age spread and delay the time at which the maximum diagram spread is reached.

The same qualitative ordering appears at both $\Sigma_{\rm cloud}=0.1$ and $1.0\,{\rm g\,cm^{-2}}$, although the absolute timescales remain shorter for the higher-surface-density models. The consistency of the trend across the H--R diagram and all three CMD projections indicates that the mass dependence is not produced by a particular photometric transformation.

\begin{figure*}
    \centering
    \includegraphics[width=1\linewidth]
    {science-images/spread_evolution_both_sigmas.png}
    \caption{Evolution of the quartile-tail diagram spread for representative $M_{\rm cl}=300$, $3000$, and $30000\,M_\odot$ simulations at fixed $\epsilon_{\rm ff}=0.03$. Results are shown for both $\Sigma_{\rm cloud}=0.1$ and $1.0\,{\rm g\,cm^{-2}}$ and for the H--R diagram and three CMD projections used in the main analysis. Each curve represents one simulation realization rather than an ensemble average. Increasing cluster mass shifts the development and peak of the diagram spread toward later times and generally produces a larger maximum spread, reflecting the longer formation timescale of the more massive clusters. \textbf{\textcolor{red}{INCREASE FONT SIZE ON THIS FIGURE}}}
    \label{fig:mass-spread-evolution}
\end{figure*}

The individual CMD evolution underlying these spread curves is shown in Appendix~\ref{app:mass-cmd-evolution}. In the more massive simulations, the stellar population remains distributed across a broader range of isochrone ages for a longer interval, visually demonstrating the extended formation history responsible for the delayed and enhanced spread.

\textbf{\textcolor{red}{AG: The subsequent subsections could go to the Discussion section.}}
\subsection{Implications for Synthetic CMD Diagnostics}

The ensemble and sensitivity analyses establish a forward connection between cloud-scale star-formation parameters, cluster mass, intrinsic stellar age distributions, and observable diagram morphology. In the simulations considered here, lower $\epsilon_{\rm ff}$ and lower $\Sigma_{\rm cloud}$ both produce more extended star-formation histories and larger diagram spreads. The cluster-mass experiment additionally shows that, at fixed $\epsilon_{\rm ff}$ and $\Sigma_{\rm cloud}$, increasing $M_{\rm cl}$ lengthens the absolute formation timescale, delays the peak in diagram spread, and generally allows a larger maximum spread to develop. Thus, CMD or H--R diagram width should not be interpreted as depending on a single cluster-formation parameter in isolation.

The correspondence is clearest when interpreted statistically across the simulation ensemble rather than from an individual snapshot. Stochastic differences between realizations and changes in isochrone geometry prevent the diagram width from serving as a unique one-to-one estimator of age spread at every epoch. Nevertheless, the persistence of the parameter ordering across H--R and multiple CMD projections shows that the width of the pre--main-sequence population contains measurable information about the underlying formation history.

The binary-resolution experiment further demonstrates that the absolute normalization of the CMD-spread statistic depends on observational angular resolution. Unresolved binaries increase the measured spread by an approximately additive offset while leaving the broad temporal behavior largely unchanged. For the representative cluster placed at 410 pc, the majority of binary systems remain unresolved under the adopted JWST and HST Rayleigh criteria, causing the predicted spread to lie considerably closer to the all-unresolved than to the all-resolved limiting case. The weak dependence on viewing projection suggests that, for this realization, the fraction of binaries that are unresolved is substantially more important than the particular orientation from which the cluster is observed.

The forward-modeling framework therefore provides a basis for connecting observed CMD morphology to physically motivated cluster-formation models. However, the absolute CMD spread cannot be interpreted solely as a measure of intrinsic stellar age spread. Quantitative comparison with an observed cluster should reproduce the relevant instrumental angular resolution and binary treatment in addition to accounting for extinction, photometric uncertainty, completeness, and field contamination. The full simulation ensemble establishes the systematic dependence of diagram spread on $\epsilon_{\rm ff}$ and $\Sigma_{\rm cloud}$, while the binary-resolution experiment identifies unresolved multiplicity as an additional observational contribution to the measured spread.





\section{Conclusion}
\label{sec:conclusion}

We have investigated the relationship between protocluster formation histories and the morphology of synthetic color--magnitude and Hertzsprung--Russell diagrams. By combining Nbody6 simulations of forming stellar clusters with stellar-evolution models and synthetic photometry generated using SPISEA, we have developed a forward-modeling framework that connects the physical parameters governing cluster formation to the observable distribution of pre--main-sequence stars on CMDs and H--R diagrams.

Our primary analysis considered 200 simulations of $3000\,M_\odot$ molecular clumps spanning two cloud surface densities and five star formation efficiencies per free-fall time. We quantified the vertical extent of the stellar sequences using a quartile-tail spread statistic measured over a fixed low-mass region of each diagram. Comparing these measurements with the intrinsic stellar age distributions revealed several principal results:

\begin{enumerate}

    \item \textbf{Star formation efficiency regulates intrinsic stellar age spread and diagram morphology.} Lower values of $\epsilon_{\rm ff}$ produce more extended star-formation histories and consequently larger intrinsic stellar age spreads. These populations generally exhibit broader stellar sequences in both H--R and CMD space. Conversely, clusters formed with higher $\epsilon_{\rm ff}$ assemble their stellar populations more rapidly and exhibit narrower age distributions and smaller diagram spreads.

    \item \textbf{Cloud surface density introduces
    an additional systematic dependence.} At fixed $\epsilon_{\rm ff}$, simulations with $\Sigma_{\rm cloud}=0.1\,{\rm g\,cm^{-2}}$ generally develop larger intrinsic age spreads and broader stellar sequences than those with $\Sigma_{\rm cloud}=1.0\,{\rm g\,cm^{-2}}$. This reflects the longer characteristic formation timescales of the lower-density molecular clouds and demonstrates that diagram morphology retains information about the physical conditions of the parent cloud.

    % \item \textbf{The relationship between intrinsic age spread and diagram width persists across different photometric representations.} 
    % The H--R diagram and the three synthetic CMD projections exhibit consistent qualitative trends with the underlying stellar age distribution. However, the absolute value and detailed time dependence of the measured spread differ between diagram spaces because of the nonlinear transformation of pre--main-sequence evolutionary tracks into photometric coordinates. Consequently, diagram spreads measured in different filter combinations are not directly interchangeable.

    \item \textbf{Increasing cluster mass
    generally increases the maximum diagram
    spread and delays its development.} A targeted comparison of representative $300$, $3000$, and $30000\,M_\odot$ simulations at fixed $\epsilon_{\rm ff}=0.03$ indicates that more massive protoclusters develop broader stellar populations over longer absolute timescales. This behavior is consistent with the longer formation timescales of more massive molecular clumps at fixed cloud surface density and star formation efficiency. Because this experiment uses one realization per mass and surface density, the result is interpreted as a qualitative sensitivity test rather than an ensemble constraint on the mass dependence.

    \item \textbf{Unresolved binary systems increase the measured CMD spread without substantially altering its overall temporal evolution.} Combining the light of binary components shifts unresolved systems toward brighter magnitudes and introduces an approximately additive contribution to the quartile-tail spread statistic. For a representative simulated cluster at 410 pc, approximately 70--80\% of current binary systems remain unresolved under the adopted JWST and HST Rayleigh criteria. The resulting spread curves lie substantially closer to the all-unresolved limit than to the all-resolved limit, while showing little dependence on the viewing projection.

\end{enumerate}

Taken together, these results establish that the vertical morphology of young stellar populations contains measurable information about their formation histories. Particularly, the systematic relationship between the duration of star formation, intrinsic stellar age spread, and diagram width survives the transformation from theoretical H--R space into multiple observable CMD projections. The persistence of these trends across independent simulation realizations indicates that they reflect the underlying cluster-formation physics rather than solely stochastic differences in the sampled stellar populations.

At the same time, our results demonstrate that diagram width is not a unique indicator of intrinsic stellar age spread. Different combinations of cloud surface density, star formation efficiency, and cluster mass can produce extended star-formation histories and consequently similar broadening of the stellar sequence. Furthermore, unresolved multiplicity modifies the absolute normalization of the measured CMD spread, while the geometry of the stellar evolutionary tracks determines how a given age distribution is expressed in different photometric systems. These effects must be considered when interpreting the observed width of a pre--main-sequence population.

\textcolor{red}{AG: need to shorten the Conclusion and avoid repetition of statements.}

\textcolor{red}{AG: Limitations and caveats can be added as a sub-section in Discussion section.}

Several limitations remain in applying the present results directly to observed clusters. Our synthetic photometry assumes a common distance, solar metallicity, and zero extinction, and the ensemble analysis does not include observational photometric uncertainties, completeness effects, or field-star contamination. The binary-resolution experiment accounts for the projected separations of bound binary systems but does not model the blending of unrelated sources or the full instrumental point-spread function. These simplifications allow the effects of intrinsic cluster formation and binary multiplicity to be examined separately, but more complete synthetic observations will be required for quantitative comparisons with individual observed clusters.

Future work can extend the framework by incorporating realistic extinction distributions, photometric uncertainties, crowding, and instrumental selection effects. A broader ensemble analysis across cluster mass would also allow the mass dependence identified here to be quantified statistically. Ultimately, these developments would support the comparison of observed young clusters with physically motivated synthetic populations, allowing their formation histories and underlying cloud conditions to be constrained using observable CMD morphology.

The present study provides a forward-modeling foundation for such comparisons by establishing which features of synthetic pre--main-sequence diagrams are sensitive to the duration and physical conditions of cluster formation, and by identifying unresolved multiplicity as an important additional contribution to their observed width.

%% IMPORTANT! The old "\acknowledgment" command has be depreciated. It was
%% not robust enough to handle our new dual anonymous review requirements and
%% thus been replaced with the acknowledgment environment. If you try to 
%% compile with \acknowledgment you will get an error print to the screen
%% and in the compiled pdf.
%% 
%% Also note that the akcnowlodgment environment does not support long amounts of text. If you have a lot of people and institutions to acknowledge, do not use this command. Instead, create a new \section{Acknowledgments}.
\begin{acknowledgments}

\end{acknowledgments}

%% To help institutions obtain information on the effectiveness of their 
%% telescopes the AAS Journals has created a group of keywords for telescope 
%% facilities.
%
%% Following the acknowledgments section, use the following syntax and the
%% \facility{} or \facilities{} macros to list the keywords of facilities used 
%% in the research for the paper.  Each keyword is check against the master 
%% list during copy editing.  Individual instruments can be provided in 
%% parentheses, after the keyword, but they are not verified.

\software{astropy \citep{2013A&A...558A..33A,2018AJ....156..123A}
          }

%% Appendix material should be preceded with a single \appendix command.
%% There should be a \section command for each appendix. Mark appendix
%% subsections with the same markup you use in the main body of the paper.

%% Each Appendix (indicated with \section) will be lettered A, B, C, etc.
%% The equation counter will reset when it encounters the \appendix
%% command and will number appendix equations (A1), (A2), etc. The
%% Figure and Table counter will not reset.

\appendix

\section{Cluster-Mass CMD Evolution}
\label{app:mass-cmd-evolution}

Figures~\ref{fig:mass-cmd-low-sigma} and
\ref{fig:mass-cmd-high-sigma} show the underlying CMD evolution for the
cluster-mass sensitivity experiment described in
Section~\ref{sec:mass-test}. The comparison is restricted to
$(F182M-F200W)$ versus F200W and $\epsilon_{\rm ff}=0.03$.

These diagrams provide the visual counterpart to the spread evolution shown
in Figure~\ref{fig:mass-spread-evolution}. At fixed cloud surface density,
the more massive clusters remain distributed across a wider range of
pre--main-sequence isochrones for a longer period of time. The characteristic
broadening and subsequent contraction of the sequence therefore occur later
as $M_{\rm cl}$ increases.

\begin{figure*}
    \centering
    \includegraphics[width=1\linewidth]
    {science-images/cmd_evolution_f182m_f200w_sigma0p1.png}
    \caption{Time evolution of the $(F182M-F200W)$ versus F200W CMD for representative $M_{\rm cl}=300$, $3000$, and $30000\,M_\odot$ simulations with $\epsilon_{\rm ff}=0.03$ and $\Sigma_{\rm cloud}=0.1\,{\rm g\,cm^{-2}}$. The more massive simulations form their stellar populations over longer absolute timescales, causing the CMD to remain broadly distributed across the pre--main-sequence isochrones for a correspondingly longer interval. \textbf{\textcolor{red}{INCREASE FONT SIZE ON THIS FIGURE}}}
    \label{fig:mass-cmd-low-sigma}
\end{figure*}

\begin{figure*}
    \centering
    \includegraphics[width=1\linewidth]
    {science-images/cmd_evolution_f182m_f200w_sigma1p0.png}
    \caption{Same as Figure~\ref{fig:mass-cmd-low-sigma}, but for $\Sigma_{\rm cloud}=1.0\,{\rm g\,cm^{-2}}$. The absolute formation timescales are shorter than in the lower-surface-density models, but the same ordering with cluster mass remains visible: larger $M_{\rm cl}$ delays the evolution of the CMD and the time at which its largest vertical spread is reached. \textbf{\textcolor{red}{INCREASE FONT SIZE ON THIS FIGURE}}}
    \label{fig:mass-cmd-high-sigma}
\end{figure*}

%% For this sample we use BibTeX plus aasjournals.bst to generate the
%% the bibliography. The sample631.bib file was populated from ADS. To
%% get the citations to show in the compiled file do the following:
%%
%% pdflatex sample631.tex
%% bibtext sample631
%% pdflatex sample631.tex
%% pdflatex sample631.tex

\bibliography{sources}{}
\bibliographystyle{aasjournal}

%% This command is needed to show the entire author+affiliation list when
%% the collaboration and author truncation commands are used.  It has to
%% go at the end of the manuscript.
%\allauthors

%% Include this line if you are using the \added, \replaced, \deleted
%% commands to see a summary list of all changes at the end of the article.
%\listofchanges

\end{document}

% End of file `sample631.tex'.
